% Modernised reading — fonds Alexandre Grothendieck, shelfmark 125, pages 1-8.
%
% This is an INTERPRETATION, not a transcription. Everything the critical
% apparatus carried moves into footnotes. In any dispute of substance the
% transcription governs: whoever cites Grothendieck must cite batch-01.fr.tex.
%
% Pass: Opus 5 (claude-opus-5), 2026-08-30 — first pass, unchecked against the
% pages by a human. The folder was read whole; it holds one batch, pages 1-8,
% and page 1 is a cover.
%
% Notation fixed for the whole document. The pages write hats for completion
% along Y and nothing else; the reading writes X-hat and says once what the hat
% means. The pages' "dim. coh." is cohomological dimension. The subscript a on
% the Ext of page 5 is a letter the folder never defines: it is kept, and the
% reading says it is undefined rather than guessing. The ambient projective
% space is P throughout, the ambient variety X of dimension r, the closed
% subscheme Y of dimension m.
%
% Substantive departures from the pages, each footnoted where it occurs:
% — page 2's boxes are given as a single chain of equivalences, which is what
%   the vertical double arrows on the sheet assert; the sheet's left-hand and
%   right-hand columns are read as statement and reduction-to-O(n), which is
%   what the arrows between them say.
% — page 2's first display is "H^p_Y(X,F) = 0 pour p > [illegible]". The bound
%   is not on the page and the reading does not supply one.
% — the folder states its duality as an equality of Ext groups; the reading
%   states it as a perfect pairing, which is what the crosses of page 5 mean
%   and what makes the statements on page 2 follow.
% — page 7's item 2 names "Lefschetz -- [name]". The reading gives the name as
%   doubtful, with the strokes described, and does not build on it.
%
% What the folder announces and never establishes: the whole of the six-item
% programme of pages 7 and 8 — it is a list of places where formal schemes are
% needed, not a proof of anything, and the folder stops two-thirds down page 8.
% Item 6, on the Greenberg-Néron types, is one sentence and has no sequel here.
\documentclass[11pt,a4paper]{article}
% Le préambule commun à toutes les transcriptions du fonds Grothendieck.
%
% Il tient en une idée : **l'incertitude se compose, elle ne se résout pas.**
% Une transcription qui lisse un mot douteux a détruit la seule chose pour
% laquelle on la faisait — un lecteur doit pouvoir savoir, à chaque endroit,
% ce qui était sur le feuillet et ce qui a été ajouté. Les six macros qui
% suivent sont donc l'appareil critique, et non de la décoration.
%
% Les mêmes commandes sont reconnues par `scripts/render.mjs`, qui produit la
% vue de lecture du site. Une macro ajoutée ici doit l'être aussi là-bas,
% faute de quoi le rendu échouera bruyamment — ce qui est voulu : un rendu
% silencieusement faux est pire qu'un rendu absent.

\ProvidesPackage{grothendieck}[2026/08/08 Transcriptions du fonds Grothendieck]

\usepackage{amsmath,amssymb,amsthm}
\usepackage{mathrsfs}
\usepackage{stmaryrd}
% Les diagrammes commutatifs sont partout dans ce fonds : c'est la langue
% même de la géométrie algébrique de Grothendieck, et une transcription qui
% les rendrait en prose ne transcrirait rien.
\usepackage{tikz-cd}
% Français par défaut — c'est la langue des feuillets — mais l'anglais est
% chargé aussi : la traduction et le résumé sont des documents anglais, et la
% typographie française (espaces avant les deux-points) y serait une faute.
% Ces documents basculent par \selectlanguage{english} en tête de corps.
\usepackage[english,french]{babel}
\usepackage{fontspec}
\usepackage{geometry}
\geometry{a4paper,margin=25mm,marginparwidth=28mm}
\usepackage{marginnote}
\usepackage{xcolor}
% ulem, not soul. `soul`'s \st and \ul are built on a character-by-character
% scan that XeLaTeX's font machinery breaks: the document fails with an
% undefined control sequence rather than degrading. ulem does the same two jobs
% and is Unicode-engine safe. `normalem` stops it hijacking \emph, which the
% transcriptions use for Grothendieck's own emphasis.
\usepackage[normalem]{ulem}
\usepackage{hyperref}
\hypersetup{colorlinks=true,linkcolor=black,urlcolor=black,pdfborder={0 0 0}}

% --- Métadonnées du lot -----------------------------------------------------
% Reprises telles quelles par le script de rendu, qui les affiche en tête de
% la vue de lecture. Elles fixent ce que le fichier prétend couvrir : sans
% elles, une transcription flotte sans point d'attache dans un fonds de seize
% mille feuillets.

\newcommand{\folder}[1]{\gdef\grothFolder{#1}}
\newcommand{\batch}[1]{\gdef\grothBatch{#1}}
\newcommand{\leaves}[2]{\gdef\grothFirst{#1}\gdef\grothLast{#2}}
\newcommand{\foldertitle}[1]{\gdef\grothFoldertitle{#1}}
\gdef\grothFolder{?}\gdef\grothBatch{?}\gdef\grothFirst{?}\gdef\grothLast{?}\gdef\grothFoldertitle{}

% --- Le résumé --------------------------------------------------------------
%
% Il ouvre la lecture modernisée, sous le titre « Résumé », et il est composé
% en petit. Ce n'est pas un ornement : c'est le seul passage du document qui
% ne soit pas des mathématiques, et un lecteur doit pouvoir le reconnaître
% avant de l'avoir lu — puis le sauter s'il connaît déjà le sujet, ou s'y
% arrêter s'il ne le connaît pas. Le filet qui le suit marque la reprise du
% corps.

\newenvironment{resume}
  {\small\setlength{\parskip}{0.45em}}
  {\par\medskip\hrule\medskip}

% --- Mots-clés ---------------------------------------------------------------
%
% En anglais, en clôture du résumé de la lecture modernisée : le vocabulaire
% moderne sous lequel chercher ce que les feuillets construisent. Le manifeste
% du site (`npm run manifest`) les extrait pour étiqueter la cote entière —
% cette ligne est la source unique des tags, il n'y a pas de fichier à côté.
\newcommand{\keywords}[1]{\par\smallskip\noindent
  {\footnotesize\textsc{Keywords} --- \textit{#1}}\par}

% --- L'appareil critique ----------------------------------------------------

% Le numéro de feuillet, en marge. C'est l'ancre qui permet de régler un
% désaccord sur une formule en regardant *un* feuillet plutôt qu'une chemise
% de six cents. Le site s'en sert aussi pour tourner les pages du fac-similé
% quand on fait défiler la transcription.
\newcommand{\leaf}[1]{%
  \marginnote{\footnotesize\color{gray!70}\textsf{#1}}%
  \label{leaf:#1}\ignorespaces}

% Illisible : on ne devine pas. Un mot inventé qui se lit comme les autres est
% le pire résultat possible.
\newcommand{\ill}{\textcolor{red!60!black}{[\,\ldots\,]}}

% Lecture douteuse : le mot est donné, et signalé comme incertain.
\newcommand{\uncertain}[1]{\uline{#1}}

% Ajout de l'éditeur — une lettre manquante, un mot sous-entendu.
\newcommand{\add}[1]{\textcolor{blue!50!black}{[#1]}}

% Biffé par Grothendieck lui-même. Ce qu'il a rayé fait partie du document :
% une rature dit souvent où la pensée a changé de direction.
\newcommand{\struck}[1]{\sout{#1}}

% Note du transcripteur, sur l'état matériel du feuillet.
\newcommand{\note}[1]{\par\smallskip{\small\color{gray!60!black}\textit{[#1]}}\par\smallskip}

% Note marginale de Grothendieck, distincte du corps du texte.
\newcommand{\marginal}[1]{\par\smallskip\noindent\hspace{1em}%
  {\small\color{gray!60!black}#1}\par\smallskip}

% --- Titre ------------------------------------------------------------------

\AtBeginDocument{%
  \begin{center}
    {\large\bfseries Fonds Alexandre Grothendieck --- cote n\textsuperscript{o}~\grothFolder}\\[2pt]
    {\itshape\grothFoldertitle}\\[4pt]
    {\small feuillets \grothFirst--\grothLast\ (lot \grothBatch)}\\[2pt]
    {\footnotesize\color{gray!60!black}Transcription établie d'après le fac-similé numérisé
     par l'Université de Montpellier.\\
     Les lectures incertaines sont soulignées, les illisibles notées [\,\ldots\,],
     les ajouts entre crochets.}
  \end{center}
  \vspace{1em}}

\endinput


\folder{125}
\pages{1}{8}
\dating{[avant 1970]}
\watermark{Édition de démonstration\\interprétation personnelle de l'œuvre}
\foldertitle{Dualité projective : notes manuscrites (s.d.) — lecture modernisée du dossier entier}

\begin{document}

\section*{Résumé}

\begin{resume}
Ce qu'on voit d'une figure quand on en retire un morceau, et ce qu'on en voit
quand on s'approche indéfiniment de ce morceau : le dossier dit que ce sont
deux façons de lire une même chose.

Soit une figure algébrique $X$ — une surface, une variété de dimension
quelconque — et dedans un morceau plus petit $Y$. Il y a deux manières de
regarder $Y$ depuis $X$. La première est de retirer $Y$ et d'examiner ce qui
reste, $X - Y$ : les trous du complémentaire renseignent sur la place que $Y$
occupait, comme un moule renseigne sur l'objet moulé. La seconde est de rester
sur $Y$ mais de ne pas s'y borner : on regarde $Y$ avec un voisinage
infiniment mince autour, ce qu'on appelle le \emph{complété formel} de $X$ le
long de $Y$ — non pas $Y$ seul, mais $Y$ plus la trace de la façon dont il est
posé dans $X$.

Ces deux regards ne sont pas indépendants. Il y a entre eux une dualité, du
même genre que celle qui, en algèbre linéaire, apparie un espace vectoriel à
l'espace des formes linéaires sur lui : à chaque énoncé d'annulation d'un côté
répond un énoncé de comparaison de l'autre, et l'un est exactement aussi fort
que l'autre. Le feuillet 2 du dossier est un tableau de six énoncés reliés par
ces équivalences, avec en marge, à chaque étage, le nom de ce qui les
justifie ; et au milieu du tableau, entouré d'un trait, le mot « dualité ».
C'est la page la plus dense du dossier et elle ne comporte pas une phrase :
c'est une carte, pas une rédaction.

Ce que la carte affirme est ceci. Dire que le complémentaire $X - Y$ n'a pas de
trous au-delà d'un certain degré, dire que la cohomologie de $X$ ne voit pas la
différence entre $X$ et le voisinage infiniment mince de $Y$, et dire qu'une
certaine comparaison de groupes d'extensions est bijective jusqu'à un
degré complémentaire — ce sont trois formulations du même fait. Le degré
critique est $r - m$, différence entre la dimension de la figure et celle du
morceau : c'est la codimension, et elle mesure combien de degrés de liberté on
perd en se restreignant à $Y$.

Le dossier est un dossier de travail et ne prouve rien de bout en bout. Deux
de ses feuillets sont d'ailleurs des pages d'un tapuscrit qu'il a annotées :
elles ne sont pas là par hasard, elles portent exactement la machine — le
calcul de la cohomologie d'un complémentaire par le complexe de Koszul d'une
famille d'équations — dont les feuillets manuscrits se servent sans la
redémontrer.

Les deux derniers feuillets sont d'un autre ordre et sont sans doute ce que le
dossier a de plus intéressant. Ils ne font pas de mathématiques : ils dressent
la liste, en six points numérotés, des endroits où l'on a besoin de la notion
de complété formel. Théorèmes généraux sur les morphismes propres ; théorie
des déformations ; la relation entre une variété et son voisinage infiniment
mince le long d'une section hyperplane, qui est le cadre des théorèmes de type
Lefschetz ; le formalisme de dualité lui-même ; les groupes formels ; et enfin
les types de Greenberg--Néron. C'est un inventaire de motifs, écrit à l'usage
de personne, et l'ordre dans lequel il est mis au propre n'est pas celui dans
lequel il a été écrit : le point 3 est écrit avant le point 2, une flèche le
rappelle vers le haut, et le 2 est entouré.

\keywords{local cohomology, formal completion, formal scheme, projective duality, local duality, Serre duality, Ext sheaves, Koszul complex, Čech cohomology, cohomological dimension, Lefschetz theorem, theorem on formal functions, Greenberg functor, Néron model, coherent duality}
\end{resume}

\pagerange{1}{8}
\section*{Le dossier, sa charpente et ses notations}

Huit feuillets, dont un de garde. Le feuillet 1 ne porte que le titre, de sa
main : \emph{Dualité locale et dualité \ldots\ projective}, un mot étant biffé
d'un pâté d'encre au milieu.\footnote{L'inventaire de Montpellier retient
seulement « Dualité projective ». La chemise porte donc un titre plus long que
le catalogue, et c'est le mot \emph{locale} qui manque au catalogue — celui
des deux qui dit le mieux ce que le dossier fait.}

Les sept feuillets qui suivent sont deux documents entrelacés. Les feuillets
2, 4, 5, 7 et 8 sont de sa main. Les feuillets 3 et 6 sont deux pages d'un
tapuscrit, numérotées III-18 et III-5, annotées à la main et surlignées au
feutre jaune.\footnote{On les lit, contrairement au tapuscrit de Bourbaki du
dossier 44 que l'édition passe. La raison est que le feuillet III-5 démontre
$H^{p}(\mathfrak{U}, F) \simeq H^{p+1}((f), M)$ pour $\mathfrak{U}$ le
complémentaire de $V(f)$ : c'est précisément la machine de cohomologie locale
dont les feuillets manuscrits se servent. Ce n'est pas du papier trouvé, c'est
la pièce du dossier. La main qui les annote est ronde, posée et droite, et ne
ressemble pas à l'écriture rapide des autres feuillets ; l'édition l'enregistre
et ne décide pas à qui elle est.}

Le dossier se lit en quatre stations :

\begin{itemize}
\item la carte des équivalences, qui est tout le contenu mathématique
  (feuillet 2) ;
\item la machine dont elle se sert, portée par les deux feuillets
  dactylographiés (feuillets 3 et 6) ;
\item le dictionnaire de la dualité, écrit sans un mot de prose, et les suites
  exactes appariées (feuillets 4 et 5) ;
\item le programme en six points (feuillets 7 et 8).
\end{itemize}

Les notations sont fixées ici. $P = \mathbf{P}^{n}$ est un espace projectif sur
un corps $k$, $X \subset P$ une sous-variété de dimension $r$, et
$Y \subset X$ un fermé de dimension $m$, avec $r > m > 0$. On note
$H^{*}_{Y}(X, F)$ la \emph{cohomologie locale} à support dans $Y$ — ce qui,
dans $X$, ne se voit que le long de $Y$ — et $\widehat{X}$ le \emph{complété
formel} de $X$ le long de $Y$. Le chapeau désigne partout cette complétion, et
seulement elle : $\widehat{F}$ est le complété du faisceau $F$,
$\widehat{\Omega}^{r}$ celui des formes de degré maximal. Enfin $J$ est
l'idéal de $Y$ et $Y_{n}$ le sous-schéma défini par $J^{n}$, de sorte que
$\widehat{X}$ est la limite des $Y_{n}$.

Un mot sur le nombre $r - m$, qui gouverne tout le dossier. C'est la
\emph{codimension} de $Y$ dans $X$. Elle apparaît de deux côtés : comme borne
supérieure pour les degrés où la comparaison à $\widehat{X}$ est bijective, et
comme la dimension cohomologique du complémentaire $P - X$. Les deux emplois
sont les deux faces de la dualité, et c'est ce que le feuillet 2 organise.

\pagerange{2}{2}
\section*{La carte des équivalences (feuillet 2)}

Le feuillet 2 n'est pas une rédaction : c'est un tableau. Six énoncés encadrés,
reliés verticalement par des doubles flèches, chacune annotée en marge de ce
qui la justifie ; et, dans la colonne de droite, à chaque étage, la réduction
du cas général au cas $F = \mathcal{O}(n)$ avec $n$ grand.\footnote{Cette
colonne de droite est la partie technique du feuillet et la moins lisible : ce
sont des encadrés courts, écrits vite. Ce qu'ils disent est qu'il suffit de
traiter les faisceaux $\mathcal{O}(n)$ pour $n$ grand, tout faisceau cohérent
étant quotient d'une somme de tels faisceaux.}

Le point de départ est la suite exacte de cohomologie locale, qu'il marque
d'une astérisque et qui est le seul ingrédient formel du tableau :
\[
  \cdots \to H^{p}_{Y}(X, F) \to H^{p}(X, F) \to H^{p}(X - Y, F) \to
  H^{p+1}_{Y}(X, F) \to \cdots
\]
Elle dit exactement que ce qu'on perd en retirant $Y$ est ce qui ne se voyait
que le long de $Y$ — et c'est elle qui fait passer d'un étage à l'autre de tout
ce qui suit.

Les trois premiers étages sont les suivants.

\textbf{Premier étage.}
\[
  H^{p}(X - Y, F) = 0 \qquad \text{pour } p \geqslant m .
\]
Le complémentaire n'a pas de cohomologie au-delà du degré $m$.

\textbf{Deuxième étage.}
\[
  H^{p}_{Y}(X, F) \longrightarrow H^{p}(X, F)
  \quad \text{bijectif pour } p > m, \ \text{surjectif pour } p = m .
\]
C'est le premier étage lu à travers la suite exacte : tuer la cohomologie du
complémentaire, c'est identifier la cohomologie locale à la cohomologie
totale.\footnote{Sous l'encadré de droite correspondant, il note :
« Le cas $p = m$ peut être omis, car $H^{m}(X, \mathcal{O}(-n)) = 0$ pour $n$
grand ». C'est l'annulation de Serre en degré $< \dim X$ pour un fibré très
négatif, et elle dispense de traiter séparément le cas limite.}

\textbf{Troisième étage}, relié au deuxième par une flèche portant,
entourée d'un trait, le seul mot du feuillet : \emph{dualité}.
\[
  \mathrm{Ext}^{q}_{\mathcal{O}_{\widehat{X}}}\!\left(
  \widehat{F}, \widehat{\Omega}^{r}_{X/k}\right)
  \;\longleftarrow\;
  \mathrm{Ext}^{q}_{\mathcal{O}_{X}}\!\left(F, \Omega^{r}_{X/k}\right)
  \quad \text{bijectif pour } q < r - m, \ \text{injectif pour } q = r-m .
\]

C'est le pivot. Les deux énoncés — l'annulation de la cohomologie de $X - Y$
au-delà du degré $m$, et la bijectivité de la comparaison des $\mathrm{Ext}$ à
valeurs dans $\Omega^{r}$ en dessous du degré $r - m$ — sont le \emph{même}
énoncé, lu des deux côtés de la dualité de Serre. Les bornes se répondent
comme il faut : $p \geqslant m$ d'un côté, $q < r - m$ de l'autre, et
$p + q = r$.\footnote{Le dossier écrit la dualité comme une identification de
groupes d'$\mathrm{Ext}$. Ce qu'elle est, et ce que les croix du feuillet 5
disent explicitement, est un \emph{accouplement parfait} : $H^{p}$ et
$\mathrm{Ext}^{r-p}(-, \Omega^{r})$ sont en dualité. On l'énonce ici sous
cette forme, parce que c'est elle qui rend visible pourquoi les bornes se
correspondent.}

Les deux étages suivants descendent du général au maniable :
\[
  H^{q}(X, \Omega(n)) \longrightarrow H^{q}(\widehat{X}, \widehat{\mathcal{O}}(n)),
  \qquad
  H^{q}(X, F) \longrightarrow H^{q}(\widehat{X}, \widehat{F}),
\]
tous deux bijectifs pour $q < r-m$ et injectifs pour $q = r - m$, le second
pour $F$ localement libre.\footnote{C'est la forme sous laquelle l'énoncé est
aujourd'hui le plus reconnaissable : la cohomologie de $X$ et celle de son
complété formel le long de $Y$ coïncident en dessous de la codimension. La
première formule du feuillet, « $H^{p}_{Y}(X,F) = 0$ pour $p > \ldots$ », a sa
borne illisible ; l'édition ne la supplée pas.}

Le bas du feuillet porte deux notes de sa main, l'une entourée d'un trait,
l'autre en diagonale. La première ajoute que
\[
  H^{i}(\widehat{X}, \widehat{\mathcal{O}}(n)) = 0
  \qquad \text{pour } n \text{ grand}, \ 0 < i < r-m ,
\]
qui est l'annulation dont les étages précédents ont besoin. La seconde, sur
cinq ou six lignes, n'est lisible que par fragments.

\pagerange{3}{3}
\section*{Un lemme de recollement, en tapuscrit annoté (feuillet 3)}

Le feuillet 3, numéroté III-18, est une page d'un tapuscrit. Ce qu'elle
démontre est un fait local que la théorie utilise constamment : si un
homomorphisme de faisceaux cohérents induit un isomorphisme sur les germes en
un point, il en induit un sur tout un voisinage.

La page construit un homomorphisme $u : A^{n} \to B$ entre faisceaux
quasi-cohérents sur $Y$, à partir d'une base
$(d_{i} \otimes 1)_{1 \leqslant i \leqslant n}$ de la fibre, relevée en des
sections globales $b_{i}$.\footnote{La formule $u : A^{n} \to B$ est encadrée
au crayon sur le feuillet ; c'est la seule chose que la main annotante y
souligne comme un énoncé.} Elle conclut :

\begin{quote}
il existe un voisinage ouvert $V$ de $y$ tel que la restriction $u|V$ soit un
\emph{isomorphisme} $A^{n}|V \to B|V$.
\end{quote}

L'argument est en deux temps : noyau et conoyau de $u$ sont cohérents, et un
faisceau cohérent dont le germe en un point est nul est nul au voisinage de ce
point. C'est le second temps qui demande la cohérence, et la marge le dit —
c'est la seule annotation du feuillet qui ajoute une justification et non une
correction :

\begin{quote}
puisque $Y$ est noethérien et $B$ de type fini par hypothèse
\end{quote}

\pagerange{4}{5}
\section*{Le dictionnaire, et les suites appariées (feuillets 4 et 5)}

Le feuillet 4 est un demi-feuillet — le bas en est arraché — et il n'y a pas
un mot de prose dessus. C'est le dictionnaire de la dualité, écrit en cinq
lignes.

La cohomologie locale se calcule dans l'ambiant projectif, et comme une limite
inductive sur les voisinages infinitésimaux :
\[
  H^{p}_{Y}(X, F) \;\simeq\; H^{p}_{Y}(P, F) \;=\;
  \varinjlim_{n} \mathrm{Ext}^{p}_{P}\!\left(\mathcal{O}_{Y_{n}}, F\right),
\]
tandis que la cohomologie du complémentaire est la même limite prise sur les
puissances de l'idéal :
\[
  H^{p}(X - Y, F) \;=\; \varinjlim_{n} \mathrm{Ext}^{p}_{P}\!\left(J^{n}, F\right).
\]
Le dual de la première est une limite \emph{projective}, et c'est là que le
complété formel apparaît :
\[
  \varprojlim_{n} \mathrm{Ext}^{r-p}_{P}\!\left(F,
  \mathcal{O}_{Y_{n}} \otimes \Omega^{r}_{P/k}\right)
  \;\xrightarrow{\ \sim\ }\;
  \mathrm{Ext}^{r-p}_{P}\!\left(\widehat{F}, \widehat{\Omega}^{r}_{P}\right).
\]
La lecture est directe : passer au dual échange limite inductive et limite
projective, et une limite projective de voisinages infinitésimaux \emph{est}
le complété formel. C'est ce seul échange qui fait toute la dualité du
dossier.\footnote{Les deux dernières lignes du feuillet sont écrites l'une
sous l'autre, terme à terme : la suite exacte de cohomologie locale au-dessus,
son dual au-dessous. Le feuillet s'arrête sur un $\mathrm{Ext}^{r}$ isolé, là
où la déchirure commence.}

Le feuillet 5 fait la même chose pour des paires, et c'est là que
l'appariement est écrit explicitement. Il porte deux suites exactes longues
l'une sous l'autre, reliées terme à terme par des croix :

\begin{tikzcd}[column sep=small, row sep=small, nodes={font=\scriptsize}]
  \mathrm{Ext}^{i}(X, F, G) \arrow[r] \arrow[d, no head, "\times" description]
  & \mathrm{Ext}^{i}(X', F', G') \arrow[r] \arrow[d, no head, "\times" description]
  & \mathrm{Ext}^{i+1}(X, F, G) \arrow[d, no head, "\times" description] \\
  \mathrm{Ext}^{r-i}_{a}(X, G, F)
  & \mathrm{Ext}^{r-i}(X', G', F') \arrow[l]
  & \mathrm{Ext}^{r-1-i}(X, G, F) \arrow[l]
\end{tikzcd}

Les traits verticaux ne sont pas des flèches : ce sont les croix par lesquelles
il note l'accouplement de dualité entre les deux lignes, et c'est la raison
pour laquelle la dualité a été énoncée plus haut comme un accouplement et non
comme une égalité.\footnote{L'indice $a$ des $\mathrm{Ext}$ de la seconde ligne
est une petite lettre que le dossier ne définit nulle part. Elle est conservée
telle quelle : on ne devine pas ce qu'elle abrège. Noter que la ligne du bas
descend en degré d'une unité de plus au troisième terme — $r-1-i$ et non
$r-i$ — ce qui est ce que le décalage de la suite exacte impose.}

Puis, pour les versions à support :
\[
  \mathrm{Ext}^{i}_{Y}(X; F, G) = \varinjlim_{n} \mathrm{Ext}^{i}(X, F_{n}, G),
\]
duale de
\[
  \varprojlim_{n} \mathrm{Ext}^{r-i}_{a}(X, G, F_{n}) \;\simeq\;
  \mathrm{Ext}^{r-i}(\widehat{X}, \widehat{G}, \widehat{F}),
\]
et la même chose pour la paire accentuée $(X', Y')$ — mais là il s'arrête et
porte un point d'interrogation de sa main.\footnote{Le point d'interrogation
est sur l'identification $\varprojlim_{n} \mathrm{Ext}^{r-1-i}(X', G', F'_{n})
\simeq \mathrm{Ext}^{r-1-i}(\widehat{X}', \widehat{G}', \widehat{F}')$. Ce
dont il n'est pas sûr est donc le passage du cas absolu au cas relatif, ce qui
est bien l'endroit délicat : l'échange de la limite projective et du
$\mathrm{Ext}$ demande une condition de Mittag-Leffler que le cas accentué ne
donne pas gratuitement.}

Le feuillet se termine sur une variante à trois termes, où $Y \subset Z$ et où
l'on compare les supports :
\[
  \mathrm{Ext}^{i}_{Y}(X; F, G) \to \mathrm{Ext}^{i}_{Z}(X; F, G) \to
  \mathrm{Ext}^{i}_{Z-Y}(X; F, G),
\]
duale d'une suite sur les complétés le long de $Y$ et de $Z$, avec le morphisme
$\widehat{X}_{/Y} \to \widehat{X}_{/Z}$ pour la relier. Il écrit en marge
« plus général », et c'est en effet la forme dont les deux précédentes sont des
cas particuliers.

\pagerange{6}{6}
\section*{La machine : cohomologie d'un complémentaire (feuillet 6)}

Le feuillet 6, numéroté III-5, est la seconde page de tapuscrit, et c'est celle
qui justifie que le dossier les garde toutes deux. Elle démontre que la
cohomologie du complémentaire d'un lieu d'annulation se calcule par un complexe
de Koszul.

Soit $f = (f_{1}, \ldots, f_{n})$ une famille de fonctions et $\mathfrak{U}$ le
recouvrement par les ouverts $X_{f_{i}}$. Le groupe des $p$-cochaînes alternées
de Čech s'écrit comme une limite inductive, et s'identifie à un cran du
complexe de Koszul :
\[
  C^{p}(\mathfrak{U}, F) = \varinjlim_{n} C^{p}_{n}(M),
  \qquad
  C^{p}(\mathfrak{U}, F) \;\xrightarrow{\ \sim\ }\; C^{p+1}\!\left((f), M\right).
\]
Ces isomorphismes étant compatibles aux cobords, on obtient :

\textbf{Proposition.} \emph{Si $X$ est un préschéma noethérien, ou un
schéma dont l'espace de base est quasi-compact, il existe pour tout
$p \geqslant 1$ un isomorphisme canonique fonctoriel}
\[
  H^{p}(\mathfrak{U}, F) \;\xrightarrow{\ \sim\ }\; H^{p+1}\!\left((f), M\right),
\]
\emph{et une suite exacte}
\[
  0 \to H^{0}((f), M) \to M \to H^{0}(\mathfrak{U}, F) \to H^{1}((f), M) \to 0 .
\]

Le mot « préschéma » de cet énoncé est une addition manuscrite portée
au-dessus de la ligne : le tapuscrit disait « si $X$ est noethérien », la main
corrige en « si $X$ est un préschéma noethérien ». C'est la correction d'un
énoncé et non d'une coquille — sans elle, l'hypothèse porte sur le mauvais
objet.

Le décalage d'un degré entre $H^{p}$ et $H^{p+1}$ est ce qui rend cette page
utile au reste du dossier : c'est lui qui transporte les annulations sur
$X - Y$ en annulations du complexe de Koszul, et réciproquement. Le corollaire
qui suit en donne le cas dégénéré : si les $f_{i}$ engendrent l'idéal unité
au-dessus de $\mathfrak{U}$, alors $H^{i}(\mathfrak{U}, G) = 0$ pour tout
$i > 0$ et tout $G$ quasi-cohérent.\footnote{Ce corollaire est
l'annulation de la cohomologie d'un ouvert affine, redémontrée par le complexe
de Koszul ; la dernière phrase du feuillet dit d'ailleurs « on redémontre ainsi
le \ldots », et la suite est coupée par le bord de la page. Trois autres
additions manuscrites y figurent : « au-dessus de $\mathfrak{U}$ », « celle qui
résulte de », « puisque », et une abréviation « (pacts) » pour
\emph{quasi-compacts}.}

\pagerange{7}{8}
\section*{Le programme en six points (feuillets 7 et 8)}

Le haut du feuillet 7 achève encore de la mathématique — un encadré donnant,
pour $n$ grand,
\[
  H^{i}\!\left(\widehat{P}, \Omega^{r}_{P}(-n)\right) = 0
  \quad \text{si } 0 < i < r-m,
  \qquad
  H^{r}\!\left(P, \Omega^{r}(-n)\right) \;\xrightarrow{\ \sim\ }\;
  H^{r}\!\left(\widehat{P}, \widehat{\Omega}^{r}(-n)\right),
\]
avec en marge l'équivalence qui compte : \emph{cela revient à dire que
$P - X$ est de dimension cohomologique $r - m$}.\footnote{C'est la
reformulation la plus économique de tout le feuillet 2 : toutes les
annulations et toutes les comparaisons de ce tableau tiennent dans une seule
assertion sur la dimension cohomologique du complémentaire de $X$ dans
l'espace projectif ambiant.}

Puis le dossier change de nature. Le reste des feuillets 7 et 8 ne démontre
rien : c'est une liste.

\begin{quote}
Les schémas formels s'introduisent en géométrie. Je les groupe de la façon
suivante :
\end{quote}

La phrase est écrite vite et deux de ses mots sont douteux : le complément de
« s'introduisent en » et le verbe qui ouvre la seconde proposition. Le sens,
lui, ne l'est pas — ce qui suit est un inventaire des emplois de la notion.

\begin{enumerate}
\item[1)] \textbf{Techniquement}, pour énoncer et prouver des théorèmes
  généraux comme les deux théorèmes fondamentaux des morphismes
  propres.\footnote{Ce sont le théorème de finitude et le théorème des
  fonctions formelles, qui est précisément l'énoncé de comparaison entre la
  cohomologie d'une fibre et celle du complété formel le long de cette fibre.
  Nommer ce dernier « fondamental » à côté de la finitude est ce que fera
  EGA III.}
\item[3)] En \textbf{théorie formelle des Modules}, pour une notion de
  familles \ldots\footnote{L'ordre est celui du feuillet : l'item 3 est écrit
  avant l'item 2, et une flèche partant de la marge gauche le rappelle vers le
  haut. La fin de la phrase est illisible. C'est la théorie des déformations,
  où le complété formel est ce sur quoi vivent les déformations
  infinitésimales.}
\item[2)] En \textbf{théorie de Lefschetz}, la relation entre une variété
  (projective, ou locale) et son complété formel le long d'une « section
  hyperplane ».\footnote{Le nom qui suit \emph{Lefschetz} sur le feuillet,
  après un tiret, est écrit d'une main rapide ; les lettres sont compatibles
  avec \emph{Grauert}, que la matière rendrait naturel, mais le feuillet ne le
  prouve pas et l'édition ne construit rien dessus. Le \emph{2} de cet item est
  entouré, ce qui est vraisemblablement la marque de l'ordre où il voulait le
  mettre.}
\item[4)] En \textbf{formalisme de dualité}, pour les énoncés où les groupes de
  cohomologie locale $H^{p}_{Y}(X, F)$ s'expriment par des $\mathrm{Ext}$ ou
  des $H^{p}$ sur $\widehat{X}$, complété formel le long de $Y$. C'est le sujet
  des six premiers feuillets, ici réduit à une ligne de la liste.
\item[5)] \textbf{Théorie des groupes formels}, où les théorèmes \ldots\ sont
  \ldots\ des schémas formels.\footnote{La phrase est en grande partie
  illisible ; ce qu'on en tient est le sujet, non l'énoncé.}
\item[6)] Dans la théorie des types de \textbf{Greenberg}--\textbf{Néron}
  s'introduisent des schémas relatifs sur des anneaux qui ne sont pas
  nécessairement ceux des points précédents.
\end{enumerate}

Le dernier point est le seul que le dossier introduise par « De plus,
notons », comme un ajout après coup, et les deux noms y sont soulignés de sa
main. Il désigne le foncteur de Greenberg, qui transporte un schéma sur un
anneau de valuation discrète complet vers un schéma sur le corps résiduel, et
les modèles de Néron, où l'objet à construire est déterminé par une propriété
universelle sur la base et non par des équations. Ce que les deux ont de commun
avec le reste de la liste, c'est que l'anneau de base y est complet — donc
formel.\footnote{Le feuillet s'arrête là. Les deux tiers inférieurs sont
blancs, et le programme n'a pas de suite dans le dossier : ni développement,
ni renvoi, ni date.}

\end{document}
